\section{From equivariant class-group characters to strict Selmer classes}
\label{sp-section}

We prove assertion~\ref{sc-sha}, the uniform comparison between
class-group characters in the cyclotomic tower and the strict
Selmer group.  The argument gives all the restriction,
continuity, local-condition and twist calculations explicitly.
Its classical input is Hilbert class field reciprocity:
for a totally imaginary number field $K$, every continuous
finite-image abelian character of $G_K$ which vanishes on all
finite inertia groups factors uniquely through
\[
                  \operatorname{Cl}(K)
                    \simeq\operatorname{Gal}(H_K/K),
                  \tag{SP1}
\]
where $H_K$ is the maximal unramified abelian extension.
The Artin map sends a prime class to its arithmetic Frobenius.
We use this theorem in its usual form; we do not give a proof
of global class field theory\cite[Chapter V, \S3,
Theorems 3.5--3.6 and Example 3.9]{milne-cft}.

Fix a totally imaginary field $L$, a prime $\ell$, and $E\ge1$.
Let $c$ be the largest integer such that
$\mu_{\ell^c}\subset L$, and suppose $c\ge1$ and $c\ge2$
if $\ell=2$.
Let $\chi:G_L\to\mathbb Z_\ell^\times$ be the cyclotomic
character, and let
\[
 A=\mathbb Q_\ell/\mathbb Z_\ell,\qquad
 \rho(g)=\chi(g)^{-E},\qquad
 L_\infty=L(\mu_{\ell^\infty}),\qquad
 G_\infty=\operatorname{Gal}(\overline L/L_\infty).
 \tag{SP2}
\]
Thus $A$ with this action is the negative Tate twist
$\mathbb Q_\ell/\mathbb Z_\ell(-E)$.
All cochains and characters below are continuous, and $A$
has the discrete topology.

The cyclotomic calculation in Section~\ref{sl-section} gives
\[
        \chi(G_L)=1+\ell^c\mathbb Z_\ell,\qquad
        G_\infty=\ker\chi.
        \tag{SP3}
\]
In particular, with $\kappa=1+\ell^c$, the quotient
$G_L/G_\infty$ is topologically generated by an element
with cyclotomic value $\kappa$.
For $\nu\ge c$, put $K_\nu=L(\mu_{\ell^\nu})$.
Then $[K_\nu:L]=\ell^{\nu-c}$, and there is a generator
$\gamma_\nu$ with
$\gamma_\nu(\zeta)=\zeta^\kappa$ for a primitive
$\ell^\nu$-th root $\zeta$.  Thus these fields and generators
satisfy exactly (SC3).

\begin{lemma}\label{sp-restriction}
Restriction induces an injective map
\[
 H^1(G_L,A)\longrightarrow
                \operatorname{Hom}_{\rm cont}(G_\infty,A).
 \tag{SP4}
\]
\end{lemma}
\begin{proof}
The action on $G_\infty$ is trivial, so a cocycle restricts
to a homomorphism and every coboundary restricts to zero.
Suppose a cocycle $z$ restricts to zero.  The identity
\[
                     z(gh)=z(g)+\rho(g)z(h)
                     \tag{SP5}
\]
shows that it is constant on cosets of $G_\infty$.
It descends continuously to the quotient
$\Gamma=G_L/G_\infty$.  Continuity follows from the quotient
topology, since the quotient map is open.

Let $\tau$ be the topological generator with
$\chi(\tau)=\kappa$.  The nonzero scalar
$\kappa^{-E}-1\in\mathbb Z_\ell$ acts surjectively on $A$:
it is a unit times a power of $\ell$, and $A$ is
$\ell$-divisible.  Choose $b\in A$ such that
\[
                        (\kappa^{-E}-1)b=z(\tau).
\]
Subtract the coboundary $g\mapsto(\rho(g)-1)b$.
The resulting cocycle vanishes at $\tau$, hence at every
positive and negative integral power of $\tau$ by (SP5).
Those powers are dense in $\Gamma$, and the cocycle is
continuous with Hausdorff target.  It therefore vanishes
everywhere.  The original cocycle was a coboundary,
proving injectivity.
\end{proof}

The nonzero twist is essential to this proof.
For the trivial action, $H^1(\mathbb Z_\ell,A)$ contains
the nonzero characters
$x\mapsto \ell^{-r}x\bmod\mathbb Z_\ell$.
No vanishing statement for that action is used here.

\begin{lemma}\label{sp-finite-level}
Let $S$ be a finite set of distinct elements of
$\mathcal S_\ell(L,E)$.  For all sufficiently large
$\nu\ge c$, there is an injective map of sets
\[
                    S\longrightarrow
                    T_{\gamma_\nu,\kappa^E,\ell^\nu}.
                    \tag{SP6}
\]
The sufficient value of $\nu$ may depend on $S$.
\end{lemma}
\begin{proof}
For each $x\in S$ choose a continuous cocycle
$z_x:G_L\to A$ representing it.
The image of each $z_x$ is finite, because $G_L$ is compact
and $A$ is discrete.  As $S$ is finite, one power
$\ell^r$ annihilates all values of all these cocycles.

For each place $v\mid\ell$ choose one decomposition group
$D_v\subset G_L$.  The strict local condition says that
there exists $b_{x,v}\in A$ with
\[
                 z_x(h)=(\rho(h)-1)b_{x,v}
                         \qquad(h\in D_v).
                 \tag{SP7}
\]
There are only finitely many pairs $(x,v)$.
Increase $r$ so that $\ell^r$ also annihilates every
$b_{x,v}$.  Choose any $\nu\ge\max(c,r)$ and write
$K=K_\nu$, $G_K=\operatorname{Gal}(\overline L/K)$.

For $h\in G_K$, $\chi(h)\equiv1\pmod{\ell^\nu}$.
Consequently $\rho(h)$ acts trivially on $A[\ell^r]$.
It follows from (SP5) that $z_x|_{G_K}$ is a homomorphism,
with values in $A[\ell^\nu]$.

We verify its local conditions at every prime of $K$.
For arbitrary $g,h\in G_L$, the cocycle identity gives
\[
 z_x(ghg^{-1})
       =\rho(g)z_x(h)+(1-\rho(h))z_x(g).
       \tag{SP8}
\]
Here $\rho(ghg^{-1})=\rho(h)$ because the cyclotomic
character is abelian.  Combining (SP7) and (SP8), on the
conjugate decomposition group $gD_vg^{-1}$ the cocycle
is the coboundary with value
\[
                          \rho(g)b_{x,v}-z_x(g).
\]
This value is still annihilated by $\ell^r$, uniformly
in $g$.  Its coboundary vanishes on
$G_K\cap gD_vg^{-1}$.  These intersections are the
decomposition groups of primes of $K$ above $\ell$,
with the corresponding embeddings into $\overline L$.
Thus $z_x|_{G_K}$ is zero on each such decomposition group.

At a prime of $L$ away from $\ell$, the cyclotomic character
is unramified, so the action on inertia is trivial.
The unramified local condition on $x$ therefore implies
that $z_x$ itself, not just its cohomology class, is zero
on inertia: every coboundary for a trivial action is zero.
This remains true on conjugate inertia groups by (SP8).
Inertia groups at primes of $K$ away from $\ell$ are
subgroups of those groups.  We have proved that the
homomorphism $z_x|_{G_K}$ is unramified everywhere.
The unramifiedness of the cyclotomic character away from
$\ell$ can be seen at each finite level from
$X^{\ell^a}-1$.  Its reduction has $\ell^a$ distinct roots
in the algebraic residue field.  Hensel's lemma lifts all
of them to the maximal unramified local extension.

Apply (SP1).  There is a character
$\theta_x:\operatorname{Cl}(K)\to A[\ell^\nu]$
whose pullback to $G_K$ is $z_x|_{G_K}$.
Use the isomorphism
\[
       \mathbb Z/\ell^\nu\mathbb Z
              \longrightarrow A[\ell^\nu],
       \qquad \overline u\longmapsto
                    u/\ell^\nu\bmod\mathbb Z_\ell
                    \tag{SP9}
\]
to regard its values as residues, and define
$\varphi_x(P)=\theta_x([P])$.
As principal divisors have trivial class,
\[
                       \sum_Pv_P(\alpha)\varphi_x(P)=0.
\]

Choose $g\in G_L$ inducing $\gamma_\nu$ with $\chi(g)=\kappa$,
as permitted by (SP3).
For $h\in G_K$, equation (SP8) simplifies to
\[
                         z_x(ghg^{-1})=\kappa^{-E}z_x(h),
                         \tag{SP10}
\]
because $\rho(h)-1$ annihilates $z_x(g)$.
Conjugation by $g$ takes the arithmetic Frobenius of $P$
to that of $\gamma_\nu P$ in the maximal unramified
abelian quotient.  To see that this quotient is preserved,
conjugation preserves finite abelian unramified extensions
of $K$ and hence their maximal such extension $H_K$;
the Frobenius assertion follows on residue fields.
Thus
\[
                \varphi_x(\gamma_\nu P)
                         =\kappa^{-E}\varphi_x(P),
 \quad\text{or equivalently}\quad
                \varphi_x(\gamma_\nu^{-1}P)
                         =\kappa^E\varphi_x(P).
\]
Together with the principal-divisor relation, these are
exactly (SC6), including its inverse on $\gamma_\nu$.

Finally suppose $\varphi_x=\varphi_y$.
Prime ideal classes generate $\operatorname{Cl}(K)$,
so $\theta_x=\theta_y$, and hence
$z_x|_{G_K}=z_y|_{G_K}$.
Their restrictions to $G_\infty\subset G_K$ are equal.
Lemma~\ref{sp-restriction} implies $x=y$.
This proves the injection (SP6).
\end{proof}

\begin{theorem}\label{sp-sha}
Assume that a fixed integer $B$ satisfies
\[
                   |T_{\gamma,\kappa^E,\ell^\nu}|\le B
\]
for every cyclic datum (SC3) over $L$.
Then $\mathcal S_\ell(L,E)$ is finite and
\[
              \#\mathcal S_\ell(L,E)
                    =\ell^{s_\ell(L,E)}\le B.
                    \tag{SP11}
\]
\end{theorem}
\begin{proof}
For every finite subset $S$ of $\mathcal S_\ell(L,E)$,
Lemma~\ref{sp-finite-level} and the assumed uniform bound
give $|S|\le B$.
If the group had more than $B$ elements, one could choose
a subset of $B+1$ distinct elements, a contradiction.
Thus the group is finite of order at most $B$.

Every element has $\ell$-power order.  A sufficiently
large power $\ell^a$ therefore kills the finite group.
The filtration by $\ell^i\mathcal S_\ell(L,E)$ has finite
$\mathbb F_\ell$-vector-space quotients.  The sum of their
dimensions is its $\mathbb Z_\ell$-length, and multiplying
their cardinalities gives
$\#\mathcal S_\ell(L,E)=\ell^{s_\ell(L,E)}$.
\end{proof}

This proves assertion~\ref{sc-sha}, with the exact local
conditions and uniform quantifiers in its source statement.
The separate unit calculation in Theorem~\ref{su-count}
proves assertion~\ref{sc-count}, and
Theorem~\ref{sl-enlargement} proves assertion~\ref{sc-layer}.
The homological finite-set bound in assertion~\ref{sc-few}
is proved subsequently in Theorem~\ref{sw-few}.
There is no use of generic injectivity of restriction
in a finite extension of degree divisible by $\ell$:
the injective restriction proved above is to the
infinite cyclotomic kernel, and the finite-level maps
are constructed separately for each finite set.

\begin{corollary}\label{sp-conditional-coarse}
Assume the finite-field homology comparison in
assertion~\ref{sc-few}, with its displayed hypotheses,
for every totally imaginary number field and every cyclic
datum to which that assertion applies.
Let $L$ be totally imaginary, let $E\ge1$, and let $\ell$
be prime.  Suppose $\mu_\ell\subset L$ if $\ell$ is odd,
and $\mu_4\subset L$ if $\ell=2$.
Write $n=[L:\mathbb Q]$, $D=|\operatorname{disc}L|$,
and let $S$ be the primes above $\ell$.
Then the strict Selmer group is finite and
\[
 \log\#\mathcal S_\ell(L,E)
       \le\exp\bigl((2En)^{86}\log(2D)\bigr).
 \tag{SP12}
\]
Using Poitou--Tate and local Tate duality in the forms
stated before Proposition~\ref{ht-h2-from-selmer}, one also has
\[
 \log\#H^2(G_{L,S},\mathbb Z_\ell(E+1))
       \le\exp\bigl((2En)^{88}\log(2D)\bigr),
 \tag{SP13}
\]
and this cohomology group is finite.
Thus this modular implication requires only the finite-set
comparison, which Theorem~\ref{sw-few} supplies.
The classical reciprocity and duality theorems used in this
implication remain the explicitly stated inputs.
\end{corollary}
\begin{proof}
Section~\ref{ac-section} proves the exact arithmetic complex
and orbit estimate.  Proposition~\ref{ht-arithmetic-numerical}
therefore gives (SC20) with $C=13$.
Corollary~\ref{sl-class-number} gives the class-number bound
(SC15).  Theorems~\ref{su-count}, \ref{sp-sha}, and
\ref{sl-enlargement} prove the other three comparison
assertions.  Proposition~\ref{sc-exponent} now applies and
gives $5(\max(13,9)+2)+11=86$, proving (SP12).
The assumed roots of unity imply $\mu_\ell\subset L$
also when $\ell=2$.  Proposition~\ref{ht-h2-from-selmer}
applied with $C=86$ gives $C_K=86+2=88$ and (SP13).
\end{proof}



